% =============================================================================
% decorations.tex — Ornements du mémoire (palette bleu clair)
%
% Style : ornements gravés « vectorian » (pgfornament), déclinés en bleu clair.
% Inspiré des mémoires classiques à cadre ornemental, avec un vocabulaire
% graphique propre : cadre à double filet + coins fleuronnés + bandeau répété,
% filets à fleuron pour les chapitres et les sections.
%
% Toutes les commandes publiques sont préfixées \btp...
% =============================================================================

\usepackage{tikz}
\usetikzlibrary{calc}
\usepackage{pgfornament}
\usepackage{eso-pic}
\usepackage{lettrine}

% -----------------------------------------------------------------------------
% Palette — couleur principale #3B5D9C
% Les teintes d'ornement (ciel, clair, pâle, profond) sont mélangées
% à partir de btpblue, pour que cadres, filets et fleurons restent
% dans la même famille.
% -----------------------------------------------------------------------------
\definecolor{btpblue}{HTML}{3B5D9C}           % bleu principal
\colorlet{btpsky}{btpblue!58!white}           % accent (filet intérieur, volutes)
\colorlet{btplight}{btpblue!30!white}          % filets, pieds de page
\colorlet{btppale}{btpblue!10!white}           % fonds d'encadrés et tableaux
\colorlet{btpdeep}{btpblue!58!black}           % titres, numéros de chapitre
\definecolor{btpslate}{RGB}{70, 82, 108}      % gris-bleu (texte secondaire)
\definecolor{btpink}{RGB}{26, 34, 43}          % encre du corps de texte
\definecolor{btpgold}{RGB}{168, 138, 82}      % rehaut discret, accordé au royal

% Alias de compatibilité
\colorlet{maindeco}{btpblue}
\colorlet{decoSlate}{btpslate}
\colorlet{btpaccent}{btpsky}

% -----------------------------------------------------------------------------
% Numéros des ornements vectorian retenus
% (redéfinissables pour changer tout le vocabulaire graphique d'un coup)
% -----------------------------------------------------------------------------
\newcommand{\btpOrnCorner}{63}   % coin fleuronné (page de garde)
\newcommand{\btpOrnCornerLight}{41} % coin à volute (pages courantes)
\newcommand{\btpOrnBand}{88}     % filet fin (bandeaux de cadre)
\newcommand{\btpOrnRule}{85}     % filet à fleuron central (séparateurs)
\newcommand{\btpOrnFlourish}{46} % volute horizontale (chapitres)
\newcommand{\btpOrnSmall}{80}    % petit motif (en-têtes, pieds de page)
\newcommand{\btpOrnLeaf}{75}     % fleuron feuillu (titres liminaires)

% -----------------------------------------------------------------------------
% Séparateurs horizontaux
% -----------------------------------------------------------------------------
% \btpRule[largeur][couleur] : séparateur centré à fleuron
\newcommand{\btpRule}[2][0.55\linewidth]{%
  \begingroup
  \centering
  \pgfornament[width=#1, color=#2]{\btpOrnRule}%
  \par
  \endgroup
}

% Séparateur de section : volute fine sous le titre
\newcommand{\btpSectionRule}{%
  \vspace{-0.35\baselineskip}%
  \noindent
  \textcolor{btplight}{\rule{\linewidth}{0.9pt}}%
  \par\vspace{-0.55\baselineskip}%
  \noindent
  \textcolor{btpblue}{\rule{0.22\linewidth}{2.2pt}}%
  \par\vspace{0.35\baselineskip}%
}

% Fleuron de fin de partie
\newcommand{\btpEndMark}{%
  \par\vspace{0.8em}
  \begin{center}
    \pgfornament[width=3.4cm, color=btpsky]{\btpOrnSmall}
  \end{center}
  \vspace{0.4em}
}

% Fleuron large pour les pages liminaires
\newcommand{\btpLiminaryRule}{%
  \begin{center}
    \pgfornament[width=4.6cm, color=btpblue]{\btpOrnFlourish}
  \end{center}
}

% -----------------------------------------------------------------------------
% Cadre ornemental de la page de garde
% -----------------------------------------------------------------------------
% Double filet + bandeau d'ornements répétés + coins fleuronnés.
% \btpCoverFrame  : à appeler dans l'environnement titlepage.
\newcommand{\btpCoverFrame}{%
  \AddToShipoutPictureFG*{%
    \begin{tikzpicture}[remember picture, overlay]
      % --- filets ---
      \draw[btpblue, line width=1.5pt]
        ($(current page.south west)+(1.15cm,1.15cm)$)
        rectangle
        ($(current page.north east)+(-1.15cm,-1.15cm)$);
      \draw[btpblue!70!white, line width=0.55pt]
        ($(current page.south west)+(1.42cm,1.42cm)$)
        rectangle
        ($(current page.north east)+(-1.42cm,-1.42cm)$);
      % --- bandeaux d'ornements répétés (haut / bas) ---
      \pgfornamentline[color=btpblue!62!white]
        {$(current page.north west)+(2.6cm,-1.285cm)$}
        {$(current page.north east)+(-2.6cm,-1.285cm)$}{9}{\btpOrnBand}
      \pgfornamentline[color=btpblue!62!white]
        {$(current page.south west)+(2.6cm,1.285cm)$}
        {$(current page.south east)+(-2.6cm,1.285cm)$}{9}{\btpOrnBand}
      % --- coins fleuronnés ---
      \node[anchor=north west, inner sep=0] at ($(current page.north west)+(1.18cm,-1.18cm)$)
        {\pgfornament[width=2.35cm, color=btpblue]{\btpOrnCorner}};
      \node[anchor=north east, inner sep=0] at ($(current page.north east)+(-1.18cm,-1.18cm)$)
        {\pgfornament[width=2.35cm, color=btpblue, symmetry=v]{\btpOrnCorner}};
      \node[anchor=south west, inner sep=0] at ($(current page.south west)+(1.18cm,1.18cm)$)
        {\pgfornament[width=2.35cm, color=btpblue, symmetry=h]{\btpOrnCorner}};
      \node[anchor=south east, inner sep=0] at ($(current page.south east)+(-1.18cm,1.18cm)$)
        {\pgfornament[width=2.35cm, color=btpblue, symmetry=c]{\btpOrnCorner}};
    \end{tikzpicture}%
  }%
}

% Cadre allégé pour la page de titre intérieure et les pages liminaires fortes
\newcommand{\btpLightFrame}{%
  \AddToShipoutPictureFG*{%
    \begin{tikzpicture}[remember picture, overlay]
      \draw[btpblue!38!white, line width=0.9pt]
        ($(current page.south west)+(1.3cm,1.3cm)$)
        rectangle
        ($(current page.north east)+(-1.3cm,-1.3cm)$);
      \node[anchor=north west, inner sep=0] at ($(current page.north west)+(1.33cm,-1.33cm)$)
        {\pgfornament[width=1.5cm, color=btpblue]{\btpOrnCornerLight}};
      \node[anchor=south east, inner sep=0] at ($(current page.south east)+(-1.33cm,1.33cm)$)
        {\pgfornament[width=1.5cm, color=btpblue, symmetry=c]{\btpOrnCornerLight}};
    \end{tikzpicture}%
  }%
}

% -----------------------------------------------------------------------------
% Décor optionnel des pages courantes : filet dans la marge extérieure
% Le décor principal des pages courantes est porté par l'en-tête et le pied
% de page (voir preamble.tex) ; ce filet est un supplément facultatif.
% -----------------------------------------------------------------------------
\newcommand{\btpEnableMarginRule}{%
  \AddToShipoutPictureBG{%
    \begin{tikzpicture}[remember picture, overlay]
      \draw[btplight, line width=2.2pt]
        ($(current page.north east)+(-1.15cm,-2.6cm)$) --
        ($(current page.south east)+(-1.15cm,2.6cm)$);
    \end{tikzpicture}%
  }%
}

% Conservé pour compatibilité : n'active plus le filet latéral.
\newcommand{\btpEnablePageDecor}{}

% -----------------------------------------------------------------------------
% Ouverture de chapitre
% -----------------------------------------------------------------------------
% Bloc bleu clair portant le numéro, mention verticale « CHAPITRE »,
% titre en petites capitales, filet à volute.
\newcommand{\btpChapterLabel}{Chapitre}

\newcommand{\btpChapterNumberBlock}[1]{%
  \begin{tikzpicture}[baseline=(blk.center)]
    \node[
      draw=btpblue, line width=1pt,
      fill=btppale,
      rounded corners=1.5pt,
      minimum width=2.05cm, minimum height=2.05cm,
      inner sep=0pt
    ] (blk) {\fontsize{46}{46}\selectfont\bfseries\color{btpdeep}#1};
    \node[
      rotate=90, anchor=south,
      font=\scriptsize\bfseries\scshape,
      text=btpblue,
      inner sep=3pt
    ] at (blk.west) {\btpChapterLabel};
  \end{tikzpicture}%
}

% -----------------------------------------------------------------------------
% Paragraphe d'accroche de chapitre (lettrine)
% -----------------------------------------------------------------------------
\setlength{\DefaultFindent}{2pt}
\renewcommand{\LettrineFontHook}{\color{btpblue}}

% \btpIntro{D}{ans ce chapitre, ...}
%   #1 : lettre initiale (mise en lettrine)
%   #2 : suite du paragraphe d'accroche
\newcommand{\btpIntro}[2]{%
  \begingroup
  \itshape
  \lettrine[lines=3, lhang=0.14, loversize=0.05, findent=3pt, nindent=0pt]%
    {#1}{}#2%
  \par
  \endgroup
  \vspace{0.6em}
  \noindent
  \hbox to \linewidth{\hfil\pgfornament[width=3cm, color=btplight]{\btpOrnSmall}\hfil}%
  \par\vspace{0.5em}
}

% -----------------------------------------------------------------------------
% Bandeau d'image de chapitre (optionnel)
% -----------------------------------------------------------------------------
\newcommand{\btpChapterBanner}[1]{%
  \noindent
  \begin{tikzpicture}
    \clip[rounded corners=3pt]
      (0,0) rectangle (\linewidth, 2.6cm);
    \node[anchor=south west, inner sep=0] at (0,0)
      {\includegraphics[width=\linewidth]{#1}};
    \fill[btpblue, opacity=0.12] (0,0) rectangle (\linewidth, 2.6cm);
  \end{tikzpicture}%
  \par\vspace{0.4cm}%
}

% Alias documenté dans le README : placer juste avant \chapter{...}
\newcommand{\setchapterimage}[1]{\btpChapterBanner{#1}}
